\section{Methodology}\label{sec:method}
% =============================================================================================
% Owner: methodology writer. Describes what src/synthval does (twosample/critic.py, nets.py,
% c2st.py, mmd.py, signature.py, ks.py, relative.py; pipeline.py; experiments/e3_sp500.py for
% the verdict). Formal statements live in Section 4 and are only referenced here.
% Equations defined here: eq:wgan-gp, eq:critic-stat, eq:mmd-u, eq:lr-oracle, eq:perm-pvalue,
% eq:dk-var, eq:hac-z, eq:anchored-obj, eq:anchored-stat, eq:two-critic-stat, eq:pmed,
% eq:verdict. Tables: tab:method-tests, tab:method-defaults. Figure: fig:alg-verdict.
% Numbers (referee revision): E1 sizes 0.085/0.080 (joint), 0.270/0.155 (indep) from
% tables/e1_size_sim.tex; D1/D2 sizes and power from results/e3/e3_power_summary.csv and
% e3_realvol_summary.csv (single-seed boundary intervals are Clopper--Pearson, as in
% tables/e3_power.tex); learning-rate discordant pairs 19:0 from results/e5/e5_simpower_raw.csv
% (family a, 1000 steps, B=1000, alpha=0.05); protocol dates from docs/DESIGN.md change log and
% the written/quick fields of results/{e3,e3_quick}/meta_*.json.
% =============================================================================================

This section defines the tests exactly as they are implemented in the code that produces every
number in the paper. All of them share one template: a test function is learned on a
\emph{training half} of the two samples and evaluated, frozen, on the held-out
\emph{evaluation half}. What changes from one setting to another is how the evaluation
statistic is calibrated. For independent windows we use permutation p-values, which have exact
finite-sample level. For the serially and cross-sectionally dependent windows of real data we
use a time-block split with Driscoll--Kraay standard errors. This studentisation is
asymptotically valid for a fixed critic whose evaluation data are independent of its training
data, a condition that the implemented split meets only approximately, so we measure the size of
the whole pipeline in a placebo calibrated to the data instead of assuming it. On
top of the absolute test we build an anchored relative test and a three-valued verdict with a
pre-specified power check. Formal properties are in \cref{sec:theory}; further implementation
details are in \cref{app:impl}. \Cref{tab:method-tests,tab:method-defaults} summarise the tests
and their default settings.

% ---------------------------------------------------------------------------------------------
\subsection{Windows, samples and hypotheses}\label{sec:method-setup}

\paragraph{Returns and scaling.}
The data are daily log-returns $r_{it}=\log S_{it}-\log S_{i,t-1}$ of tickers $i=1,\dots,N$.
\TRAIN{} (returns from 2018-01-03 to 2022-12-30, $\Ttr=1258$ trading days; the price of
2018-01-02 is the base of the first return) is the calibration period and \TEST{} (2023-01-03 to
2025-06-16, $\Tte=615$ days) the evaluation period. For each ticker let
$s_i$ be the standard deviation (divisor $\Ttr-1$) of its \TRAIN{} returns. Real \emph{and}
synthetic returns of ticker $i$ are divided by the same constant $s_i$ before anything else.
Given \TRAIN{}, this is a fixed bijection applied to both samples, so it leaves every per-ticker
null hypothesis unchanged; its only role is to put tickers on a common scale so that their
windows can be pooled (\cref{rem:pooling}).

\paragraph{Windows.}
A window is a block of $L$ consecutive normalised returns. Windows do not overlap (step $L$),
and an incomplete tail is dropped. The default is $L=20$, about one trading month; the
per-ticker tests of E3 and the multivariate tests of E4 use $L=5$. A univariate window lies in
$\R^{L}$. A multivariate window of $N$ assets lies in $\R^{N\times L}$ (E4: $N=30$). Each window
carries the \emph{date rank} $\tau$ of its first day on the common trading calendar: \TRAIN{}
days have ranks $0,\dots,\Ttr-1$ and \TEST{} days $\Ttr,\dots,\Ttr+\Tte-1$. Synthetic windows
simulated over the \TEST{} horizon receive the \TEST{} ranks, and windows of different tickers
that start on the same day share their rank. At $L=20$ a ticker yields 62 \TRAIN{} and 30
\TEST{} windows, so the 100-ticker panel of E1, E3 and E5 gives 6200 and 3000 windows.

\paragraph{Samples.}
The roles of the samples never change. $X_1,\dots,X_n$, with law $P$, are the real windows of
the comparison sample; on real data $X=\Rnew$, the \TEST{} windows (law $\Pnew$ in relative
statements). $Y_1,\dots,Y_m$, with law $Q$, are synthetic windows from a generator $G$ (law $\QG$
in relative statements). In E3, $G$ is fitted per ticker on \TRAIN{} and simulates one
\TEST{}-length path per ticker, started from the fitted end-of-\TRAIN{} state when the model has
one, so each ticker contributes as many synthetic as real windows. This matters: the pooled null
follows from the per-ticker nulls only when the per-ticker proportions are the same in both
samples (\cref{rem:pooling}). $Z_1,\dots,Z_k$, with law $\Ppast$, are the real past windows
$\Rpast$ (\TRAIN{}); only the relative tests use them.
In the simulations of E2, $X\sim P$ and $Y\sim Q$ have known laws and every window is simulated
from its own path after a burn-in, so the windows are independent and identically distributed
(\iid{}). On real data, and on windows cut from long simulated paths, they are not. Volatility
clustering makes the windows of a ticker serially dependent, and common market shocks make
windows of different tickers on the same dates dependent.

\paragraph{Hypotheses.}
The absolute null is $H_0: P=Q$ (\cref{def:abs-null}). It concerns the law of \emph{one}
window. Dependence \emph{between} windows is not part of $H_0$; it is a nuisance that the
inference protocol of \cref{sec:method-inference} must absorb. A window law fixes all the
dependence of the return process up to lag $L-1$ (\cref{prop:window-ident}). A window test can
therefore detect volatility clustering within a month, but not features that act only at longer
horizons.

Markets are not stationary. \TRAIN{} contains the 2020 crash and the 2022 drawdown. The average
correlation of the E4 universes fell from \TRAIN{} to \TEST{} (\cref{tab:e4-correlation}), and
so did the volatility level of the panel: the panel median of the daily standard deviation is
1.62\% in \TEST{}, against 2.11\% for the GBM volatility fitted on \TRAIN{}
(\cref{tab:e3-stylized,tab:e3-fit-diagnostics}). E1 compares the \TRAIN{} and \TEST{} windows
directly; the window tests with valid p-values do not reject, but with about 15 \TEST{} evaluation
dates they have little power for this comparison (\cref{tab:e1-nonstat,sec:emp-placebo}). A generator calibrated on \TRAIN{} inherits features of that period,
such as its volatility level, and cannot be expected to reproduce the \TEST{} law exactly, so
absolute rejections are to be expected on real data. The relative tests of
\cref{sec:method-relative} ask a weaker and more useful question: is $G$ detectably
further from $\Rnew$ than the real past $\Rpast$ is?

% ---------------------------------------------------------------------------------------------
\subsection{The Wasserstein-critic statistic}\label{sec:method-critic}

\paragraph{Split and standardisation.}
Each sample is split into a training and an evaluation half, with index sets $\Itr,\Iev$
(real) and $\Jtr,\Jev$ (synthetic) and sizes $n_t,n_e$ and $m_t,m_e$; the default is 50/50.
For independent windows the halves are drawn at random within each sample (half of the windows,
rounded, for training). For dependent windows we use the time-block split of
\cref{sec:method-inference}. The split never depends on the values of the windows. Inputs are
standardised by $A(x)=\Lambda^{-1}(x-m_A)$, where $m_A$ and the positive diagonal $\Lambda$ hold
the mean and standard deviation of the pooled training half $\{X_i\}_{i\in\Itr}\cup\{Y_j\}_{j\in\Jtr}$.
They are computed per coordinate for univariate windows and per asset channel (pooled over the
$L$ positions) for multivariate windows. The same $A$ is applied to every window, training or
evaluation, real or synthetic.

\paragraph{Training.}
The critic is $\fhat=\ghat\circ A$, where $\ghat$ is a neural network. It is trained \emph{alone},
with no generator, with the gradient-penalised Wasserstein GAN objective of
\citet{gulrajani2017improved} (see also \citealp{arjovsky2017wasserstein}):
\begin{equation}\label{eq:wgan-gp}
  \ghat\approx\argmax_{g}\;\Big\{\Ptr(g\circ A)-\Qtr(g\circ A)
  -\lgp\,\E\big(\norm{\nabla g(\tilde x)}-1\big)^2\Big\},\qquad
  \tilde x=u\,A(X)+(1-u)\,A(Y),
\end{equation}
where $\Ptr$ and $\Qtr$ are empirical means over the training halves, $X$ and $Y$ are training
windows drawn independently, $u\sim U(0,1)$ is independent of them, and $\lgp=10$. The gradient
is taken with respect to the standardised input. Each optimiser step draws $b=128$ indices with
replacement from each training half and $b$ independent values of $u$, pairs the $r$-th real
and synthetic draws to form $\tilde x_r$, and takes one Adam step \citep{kingma_2015_adam} on
the mini-batch loss
$-b^{-1}\sum_{r}\{g(A X_{i_r})-g(A Y_{j_r})\}+\lgp b^{-1}\sum_r(\norm{\nabla g(\tilde x_r)}-1)^2$.
We follow \citet{gulrajani2017improved} for the penalty weight and the Adam moment parameters
$(0.5,0.9)$, but use a learning rate of $10^{-3}$. It was chosen before any E2 experiment on the
basis of an informal check, which is not part of the reported experiments or of the repository,
in which their rate of $10^{-4}$ appeared to under-train the critic within 1000 steps and to lose
power. E5 contradicts this check in one of its two designs: on simulated windows the rate
$10^{-4}$ is more powerful against a 10\% volatility increase (0.78 against 0.40; in all 19
replications in which the two rates disagree, only $10^{-4}$ rejects) and not detectably
different against a change of persistence (\cref{sec:robust}). The default was fixed before the experiments and not
re-tuned, so the critic power reported in this paper is that of a pre-specified critic, not of an
optimised one. There is no early stopping: the weights after step 1000 are frozen.

\paragraph{Statistic.}
With $\fhat$ frozen,
\begin{equation}\label{eq:critic-stat}
  T=\Pev\fhat-\Qev\fhat=\frac1{n_e}\sum_{i\in\Iev}\fhat(X_i)-\frac1{m_e}\sum_{j\in\Jev}\fhat(Y_j).
\end{equation}
The critic is trained to score real windows high, so large values of $T$ are evidence against
$H_0$, and every critic-based test is one-sided. Given the training $\sigma$-field $\Gsig$
(training halves, split, training randomness), $T$ is unbiased for the critic discrepancy
$\Dcrit(\fhat)=P\fhat-Q\fhat$, and $\Dcrit(\fhat)\le\Lip(\ghat)\,\W(P_A,Q_A)$, where $P_A$ and
$Q_A$ are the laws of the standardised windows (\cref{prop:unbiased}). If $\ghat$ were
1-Lipschitz, $T$ would therefore estimate a lower bound on the Wasserstein-1 distance between
the standardised window laws (Kantorovich--Rubinstein duality, \cref{fact:kr-duality}). The
gradient penalty enforces the Lipschitz constraint only softly and only at the interpolates, and
WGAN critics are known to be poor estimators of $\W$ itself \citep{mallasto2019well}. We
therefore read $T$ as an \emph{approximate} lower bound and use it only as a test statistic
(\cref{rem:gp-soft}). The level guarantees of \cref{sec:method-inference} do not depend on any
Lipschitz property, on the convergence of training or on the architecture: a poor critic costs
power, never level.

Why hold data out? The plug-in distance $\W(P_n,Q_m)$ and the in-sample training objective are
biased upwards, and the bias vanishes only at rate $n^{-1/d}$ in dimension $d$
\citep{fournier2015rate}. The held-out $T$ is conditionally unbiased and has an exact null
distribution under exchangeability (\cref{rem:holdout}). The price is that only half of the
data trains the critic and only half evaluates it; \cref{prop:local-power} quantifies the loss.
The held-out-witness template is that of \citet{kirchler2020two} and \citet{kubler2022witness},
and \citet{wang2022two} split and permute for Wasserstein testing. Our contribution is not the
template but its instance, its validation and the protocols of
\cref{sec:method-inference,sec:method-relative,sec:method-decision}.

\paragraph{Architectures.}
For univariate windows $\ghat$ is a multilayer perceptron $\R^L\to\R^{128}\to\R^{128}\to\R^{128}\to\R$
with LeakyReLU activations (slope 0.2) and no normalisation layers. For multivariate windows
$(N,L)$ (E4) it is a one-dimensional convolutional network with the assets as input channels.
It has two blocks, each a convolution with 64 channels, kernel size $\min(5,L)$ and ``same''
padding, followed by a LayerNorm over the $(64,L)$ feature map and a LeakyReLU. The feature map
is then flattened and passed through a dense layer of width 64, a LeakyReLU and a linear
output. Flattening, rather than pooling over time, keeps the position of each day inside the
window. BatchNorm is never used: the gradient penalty is a per-window quantity, and BatchNorm
would couple the windows of a mini-batch. The multilayer perceptron is the network of every
test on univariate windows in this paper (E1, E2, E2-dep, E3 and E5), including the anchored test
behind the verdicts. A time-pooled convolutional critic, which averages convolutional features
over the days of a window, was examined only after the real-data results were known, and was
selected by a pilot on four E2 designs that is not part of the reported experiments; this
post-hoc ablation is reported in \cref{sec:robust} and does not replace the pre-specified
perceptron.

\begin{table}[tbp]
  \centering
  \caption{Default training settings of the learned tests (the code defaults used in every
    experiment unless stated). The anchored and two-critic relative tests use the \WC{} settings.}
  \label{tab:method-defaults}
  \small
  \begin{tabular}{@{}>{\raggedright\arraybackslash}p{0.21\linewidth}
                  >{\raggedright\arraybackslash}p{0.37\linewidth}
                  >{\raggedright\arraybackslash}p{0.36\linewidth}@{}}
    \toprule
    & \WC{} & \LRC{} and C2ST \\
    \midrule
    Network, $L$-windows & MLP $L$--128--128--128--1, LeakyReLU(0.2) & same \\
    Network, $(N,L)$-windows & CNN: 2 $\times$ [Conv1d(64, kernel $\min(5,L)$), LayerNorm, LeakyReLU], dense 64 & same \\
    Input scaling & per coordinate (per asset channel), pooled training half & same \\
    Objective & WGAN-GP \eqref{eq:wgan-gp}, $\lgp=10$ & class-balanced cross-entropy \\
    Optimiser & Adam, learning rate $10^{-3}$, $(\beta_1,\beta_2)=(0.5,0.9)$ & Adam, learning rate $10^{-3}$, $(0.9,0.999)$, weight decay $10^{-4}$ \\
    Mini-batch & 128 per sample, drawn with replacement & 128 per class \\
    Steps & 1000, no early stopping & at most 1000; early stopping (20\% validation, checked every 50 steps, patience 5, best checkpoint kept) \\
    Training / evaluation & 50/50 & 50/50 \\
    Permutations $B$ & 1000 & 1000 (accuracy: randomised ties) \\
    \bottomrule
  \end{tabular}
\end{table}

% ---------------------------------------------------------------------------------------------
\subsection{Competing tests and the likelihood-ratio oracle}\label{sec:method-competitors}

\paragraph{Likelihood-ratio critic and C2ST.}
Both use the split, standardisation and architecture of the \WC{} test, but train a classifier.
Its logit $\lhat$ (the network output) minimises the class-balanced binary cross-entropy
$\tfrac12\{\Ptr\log(1+e^{-\lhat})+\Qtr\log(1+e^{\lhat})\}$, with real windows labelled 1 and
mini-batches of 128 windows per class. The optimiser is Adam with learning rate $10^{-3}$,
default moments $(0.9,0.999)$ and weight decay $10^{-4}$, for at most 1000 steps. Early stopping
holds out a random 20\% of each training half, evaluates the balanced validation cross-entropy
every 50 steps, stops after five evaluations without improvement and restores the best
checkpoint. At $n=500$ the selected checkpoint is typically the first one (step 50), because
the classifier overfits early. A pilot of alternative settings found none with higher
power at valid size (\cref{app:impl-nets}), so we keep the defaults and record the selected step instead of tuning it.

With balanced classes, the Bayes-optimal logit is $\lstar=\log(p/q)$ (\cref{prop:bayes-logit};
\citealp{goodfellow2014generative}). The \LRC{} statistic is $\Tlogit=\Pev\lhat-\Qev\lhat$. If
$\lhat=\lstar+\text{const}$, its conditional mean is the Jeffreys divergence $\Jeff(P,Q)$
(\cref{cor:jeffreys}), so $\Tlogit$ is a learned log-likelihood-ratio statistic
\citep{cheng2022classification}. Unlike $T$ it has no lower-bound property: rescaling $\lhat$
inflates it. The C2ST statistic \citep{lopezpaz2017revisiting,kim2021classification} is the
balanced accuracy of the rule $\{\lhat>0\}$,
$\Tacc=\tfrac12+\tfrac12(\Pev h-\Qev h)$ with $h=\one\{\lhat>0\}$; its conditional mean is at
most $\tfrac12\{1+\TV(P,Q)\}$. The three learned statistics are thus affine functions of a
difference of evaluation means of one function that is fixed given $\Gsig$: the critic, the
logit or the indicator $h$. This is the only structure the inference of
\cref{sec:method-inference} uses.

\paragraph{Kernel tests.}
The MMD test \citep{gretton2012kernel} uses the unbiased estimate
\begin{equation}\label{eq:mmd-u}
  \TMMD=\frac{\sum_{i\neq i'}k(X_i,X_{i'})}{n(n-1)}+\frac{\sum_{j\neq j'}k(Y_j,Y_{j'})}{m(m-1)}
  -\frac{2}{nm}\sum_{i,j}k(X_i,Y_j),\qquad k(x,y)=e^{-\norm{x-y}^2/(2\varsigma^2)},
\end{equation}
on the flattened normalised windows, with no further standardisation. The bandwidth
$\varsigma$ is the median heuristic: the median pairwise Euclidean distance in a random subset
of at most 2000 points of the pooled sample. It depends only on the pooled sample, so it is
invariant to the labels. Samples larger than 3000 windows are subsampled without replacement;
E4 fixes the bandwidth per universe (\cref{sec:mk-design}). There is no training split: the
kernel tests use all windows.

The \SigMMD{} test \citep{chevyrev2022signature,kiraly2019kernels} maps a window
$(r_1,\dots,r_L)$ to the piecewise-linear path through the points
$(k/L,\sum_{l\le k}r_l)$, $k=0,\dots,L$: time and cumulative return. Its signature truncated at
depth 4 is computed exactly by Chen's identity \citep{lyons1998differential,chevyrev2016primer},
which gives $2+4+8+16=30$ features. Each feature is standardised with the mean and standard
deviation of the pooled sample, and $\Tsig$ is the statistic \eqref{eq:mmd-u} on these features,
with its own median-heuristic bandwidth.

\paragraph{Kolmogorov--Smirnov.}
$\TKS$ is the two-sample KS statistic between the $nL$ pooled real returns and the $mL$ pooled
synthetic returns, with a two-sided asymptotic p-value. It sees only the one-dimensional
marginal law of a daily return, so it has no power against deviations that leave this marginal
unchanged (\cref{fig:e2dep-power}). Its p-value assumes \iid{} returns. That assumption fails
when the returns of a window share a conditional variance, which makes KS oversized on GARCH
windows (\cref{tab:e2-size}), and it fails again across dependent windows. We keep KS because
practitioners use it, and because against deviations of the marginal (volatility level, tails,
jumps) it is, by a wide margin, the most powerful of the two-sample tests we compare
(\cref{sec:sim-power}).
Any honest use of a window test has to be weighed against this baseline.

\paragraph{Neyman--Pearson oracle (E2 only).}
When $P$ and $Q$ have known window densities $p,q$, the oracle tests $Y\sim P$ against
$Y\sim Q$ on the synthetic sample with
\begin{equation}\label{eq:lr-oracle}
  \TLR=\sum_{j=1}^m\log\frac{q(Y_j)}{p(Y_j)}.
\end{equation}
Its null distribution is computed from 2000 independent samples of $m$ windows drawn under $P$,
and the Monte Carlo p-value with randomised tie-breaking has exact level
\citep{dufour_2006_monte}. For i.i.d.\ returns (volatility level, tails, jumps) the window
densities are exact, so this is the most powerful test of $P$ against $Q$, and its power bounds
that of every level-$\alpha$ two-sample test (\cref{prop:oracle-bound}). For GARCH windows the
density has no closed form. We integrate over the conditional variance at the start of the
window with $M=512$ Monte Carlo draws from its law after the burn-in. The draws are shared by
the observed and the null samples, so the test keeps exact level, and $\TLR$ approaches the
exact log-likelihood ratio as $M$ grows.

% ---------------------------------------------------------------------------------------------
\subsection{Inference for independent and for dependent windows}\label{sec:method-inference}

For the learned tests, conditionally on $\Gsig$, every statistic is a difference
$\bar a-\bar b$ of evaluation means of a fixed function $g$, with $a_k=g(X_k)$, $k\in\Iev$, and
$b_k=g(Y_k)$, $k\in\Jev$. Two protocols turn this difference into a p-value.

\paragraph{Protocol for independent windows (E2).}
Pool the evaluation values, $(c_1,\dots,c_{n_e+m_e})=(a,b)$. For $b=1,\dots,B$, draw a uniform
random permutation of the pooled vector, independently of everything else, and let $T_b$ be the
statistic recomputed with the permuted labels. The p-value is
\begin{equation}\label{eq:perm-pvalue}
  \pperm=\frac{1+\#\{b\le B:\ T_b\ge T\}}{1+B},
\end{equation}
with $B=1000$; ties with $T$ count against it (up to a relative tolerance of $10^{-12}$). The
C2ST accuracy is discrete, so for it ties are broken at random:
$\pperm=(\#\{T_b>T\}+U\,(1+\#\{T_b=T\}))/(1+B)$ with $U\sim U(0,1)$, which is exactly uniform
under exchangeability. If the windows are \iid{} and the two samples are independent, the pooled
evaluation windows are exchangeable given $\Gsig$ under $H_0$. The test then has level at most
$\alpha$ in finite samples, for every critic, every $B$ and every split that does not depend on
the data (\cref{prop:exact-level}; \citealp{hemerik2018exact}). The kernel tests have no training
split. They permute the labels of the pooled sample with the kernel matrix and the bandwidth held
fixed ($B=500$), which is exact under the same conditions.

\paragraph{Why real data need another protocol.}
Two things go wrong when windows are dependent (\cref{sec:theory-dependence}).
\begin{enumerate}
  \item \emph{Variance.} The real evaluation windows of a ticker are autocorrelated through
    volatility clustering, and windows of different tickers on the same dates are correlated
    through the market. The permutation distribution ignores both, so the test is oversized
    (\cref{rem:naive-perm}).
  \item \emph{Bias.} If training and evaluation windows are close in time, the critic learns the
    realised volatility path of the training dates (``high volatility in spring 2023 means
    real''). Real evaluation windows from the same or neighbouring dates share that path;
    synthetic windows do not. Then $\E[T\mid\Gsig]>0$ under $H_0$. A random split of
    consecutive windows creates this coupling. A split by ticker in a panel creates it in full,
    because every training date is also an evaluation date (\cref{prop:ticker-bias}). No
    variance estimator can remove a bias.
\end{enumerate}
The naive protocol, a ticker split with permutation p-values, is run in E1 as a comparison.
In the simulated placebo panel it rejects a true null at nominal 5\% in all 200 replications
(\cref{tab:e1-size-sim}).

\paragraph{Protocol for dependent windows (E1, E3, E4, E5 and E2 part g).}
It has three components.

\emph{Time-block split.} The date range $[\tau_{\min},\tau_{\max}]$ of each sample is cut into
$K=10$ intervals of equal length: a window with date rank $\tau$ gets the block label
$\min\{K-1,\lfloor K(\tau-\tau_{\min})/(\tau_{\max}-\tau_{\min})\rfloor\}$. Labels are shared
across samples, so the blocks of $X$ and of $Y$ that carry the same label always fall in the
same half.
The $K$ labels are shuffled, and the training half receives them in that order up to the point
where the cumulative number of windows (over all samples) is closest to half the total. In E3,
$X$ and $Y$ both cover the 30 \TEST{} window dates. Every block then holds three dates (300
windows per sample), five blocks train and five evaluate, and each half has 15 dates and 1500
windows per sample. Training and evaluation windows never share a date, which removes the
common-date bias of \cref{prop:ticker-bias}. Two features of the implementation limit what the
theory covers. The blocks interleave in time, and there is no gap between a training block and
an adjacent evaluation block. The dependence across these boundaries is therefore not
controlled by \cref{lem:decoupling}, which needs a gap. We measure its effect instead
(\cref{tab:e2-dependent-size,tab:e1-size-sim}).

\emph{Driscoll--Kraay variance.} Let $\tau_k$ be the date rank of evaluation window $k$ and
$u_k=(\tau_k-\min_{k'}\tau_{k'})/L$ its position on the grid of window start dates, of length
$U=\max_ku_k+1$. Positions without an evaluation window (the training blocks) count as zeros.
With the date sums $\mathcal S_u=\sum_{k:\,u_k=u}(a_k-\bar a)$,
\begin{equation}\label{eq:dk-var}
  \widehat V_{\mathrm{DK}}(\bar a)=\frac{1}{n_e(n_e-1)}\Big[\sum_{u}\mathcal S_u^2
  +2\sum_{h=1}^{H}\Big(1-\frac{h}{H+1}\Big)\sum_{u}\mathcal S_u\mathcal S_{u+h}\Big].
\end{equation}
Summing within a date absorbs arbitrary cross-sectional dependence \citep{driscoll1998consistent},
and the Bartlett weights \citep{newey1987simple} handle serial dependence. The factor
$n_e/(n_e-1)$, folded into the normalisation, makes \eqref{eq:dk-var} equal to the usual \iid{}
variance estimator when $H=0$ and all dates are distinct. Lags are counted in windows of $L$
trading days, never in days. The truncation lag $H$ follows the AR(1) plug-in rule of
\citet{andrews1991heteroskedasticity} for the Bartlett kernel. The coefficient $\hat\rho$ is the
lag-one autocorrelation of $(\mathcal S_u)$, computed only over adjacent grid positions that both
carry evaluation windows (so that the holes left by the training blocks are not mistaken for low
persistence) and clipped to $[-0.97,0.97]$. Then
$\hat\eta=4\hat\rho^2/\{(1-\hat\rho)^2(1+\hat\rho)^2\}$ and
$H=\lceil1.1447(\hat\eta U)^{1/3}\rceil-1$, truncated to $[0,U-1]$. If fewer than ten adjacent
pairs are observed, the rule is applied to the sequence of observed dates. The rule of
\citet{newey1994automatic} is an E5 robustness option. The studentised statistic and its
p-value are
\begin{equation}\label{eq:hac-z}
  \zhac=\frac{\bar a-\bar b}{\sqrt{\widehat V_{\mathrm{DK}}(\bar a)+\widehat V_{\mathrm{DK}}(\bar b)}},
  \qquad \phac=1-\Phi(\zhac).
\end{equation}
The two variances add because the synthetic sample is independent of the real one. For a fixed
critic and a growing number of evaluation dates, $\zhac$ is asymptotically $N(0,1)$ under
$H_0$ when the window sequences are mixing (\cref{prop:hac}). With only 15 evaluation dates the
Driscoll--Kraay estimate is noisy and biased downwards, so some oversize remains. E1 measures the
size of exactly this pipeline and reports fixed-$b$ p-values \citep{kiefer_2005_new} as
comparison rows (\cref{rem:small-T}). At nominal 5\% the \WC{} test rejects a true null at 0.085
(the anchored relative test at 0.080) when the synthetic panel shares the factor structure of the
real one, but at 0.270 (anchored: 0.155) when the synthetic tickers are drawn independently, as
the univariate generators of E3 draw them (\cref{tab:e1-size-sim}). In that case the studentised
statistic is also shifted, not only over-dispersed (\cref{rem:small-T}), so a better variance
estimate alone would not remove the excess; \cref{sec:emp} takes this into account. Under this protocol the code also computes the permutation
p-value under the time-block split. E1 reports its size, but it is never used for inference.

\emph{Wild bootstrap for the kernel tests.} The kernel tests have no split and no witness
function, so they cannot be studentised in this way. They use the wild bootstrap for degenerate
V-statistics of \citet{chwialkowski2014wild}. Let $\mathsf K$ be the pooled kernel matrix and
$v=(\one_n/n,-\one_m/m)$. The p-value refers to the V-statistic $v^\top\mathsf Kv$, the squared
distance between the two empirical kernel mean embeddings. In replicate $b$, a stationary
Gaussian AR(1) process $(\zeta_u)$ with $\Cov(\zeta_u,\zeta_{u'})=\exp(-\abs{u-u'}/\ell_w)$ is
drawn on the pooled grid of window dates. Window $k$ receives $\zeta_{u_k}$, so windows that
start on the same date share one multiplier, in both samples. The multipliers are centred
within each sample, $\tilde\zeta$, and the replicate is $\tilde v^\top\mathsf K\tilde v$ with
$\tilde v=(\tilde\zeta_X/n,-\tilde\zeta_Y/m)$. The p-value $\pwild$ is the share of the
$B=500$ replicates at or above $v^\top\mathsf Kv$, computed as in \eqref{eq:perm-pvalue}. The
correlation length is $\ell_w=(\hat\eta U)^{1/3}$, where $\hat\eta$ is now the several-series
version of the Andrews constant \citep{andrews1991heteroskedasticity}, computed from the date
sums of the scores on the five leading kernel principal components ($\ell_w=0$, i.e.\
independent multipliers, when no dependence is detected). This plug-in is our heuristic:
\citet{chwialkowski2014wild} leave $\ell_w$ as a tuning parameter, and we give no proof of
validity in our setting (\cref{rem:wild}). The tables report $\TMMD$ and $\Tsig$ as the
unbiased statistic \eqref{eq:mmd-u}, while $\pwild$ refers to the V-statistic.

\emph{KS} is computed and reported, but flagged as invalid under this protocol.

\begin{table}[tbp]
  \centering
  \caption{The tests as implemented. ``Dependent windows'' is the protocol of real data and of
    windows cut from long simulated paths; ``DK'' is the Driscoll--Kraay studentisation
    \eqref{eq:hac-z} after the time-block split.}
  \label{tab:method-tests}
  \footnotesize
  \setlength{\tabcolsep}{4pt}
  \begin{tabular}{@{}>{\raggedright\arraybackslash}p{0.13\linewidth}
                  >{\raggedright\arraybackslash}p{0.27\linewidth}
                  >{\raggedright\arraybackslash}p{0.15\linewidth}
                  >{\raggedright\arraybackslash}p{0.18\linewidth}
                  >{\raggedright\arraybackslash}p{0.19\linewidth}@{}}
    \toprule
    Test & Statistic (target) & Learned on training half & p-value, independent windows &
    p-value, dependent windows \\
    \midrule
    \WC{} & $T=\Pev\fhat-\Qev\fhat$ (approximate lower bound on $\W(P_A,Q_A)$) & WGAN-GP critic &
      permutation, $B=1000$ & DK $\phac$ \\
    \LRC{} & $\Tlogit=\Pev\lhat-\Qev\lhat$ ($\Jeff(P,Q)$ if $\lhat=\lstar$) & classifier &
      permutation, $B=1000$ & DK $\phac$ \\
    C2ST & balanced accuracy $\Tacc$ (at most $\tfrac12\{1+\TV(P,Q)\}$) & classifier &
      permutation, randomised ties & DK $\phac$ \\
    MMD & $\TMMD$, Gaussian kernel on windows & none & permutation, $B=500$ &
      wild bootstrap, $B=500$ \\
    \SigMMD{} & $\Tsig$, Gaussian kernel on depth-4 signatures & none & permutation, $B=500$ &
      wild bootstrap, $B=500$ \\
    KS & $\TKS$ on pooled returns (one-dimensional marginal) & none & asymptotic (needs \iid{}
      returns) & reported, marked invalid \\
    \addlinespace
    Anchored relative & $\Danc=\Ppast^{\evl}\fhat-\QG^{\evl}\fhat$ & one critic, $\Rnew$ vs pool &
      permutation of $G$ / $\Rpast$ labels & DK (verdict p-value) \\
    Two-critic relative & $\Dtwo$, \eqref{eq:two-critic-stat} & two critics &
      paired bootstrap, $B=500$ & DK (secondary) \\
    \addlinespace
    LR oracle (E2) & $\TLR$, \eqref{eq:lr-oracle} & none (true densities) &
      Monte Carlo, 2000 null samples & -- \\
    \bottomrule
  \end{tabular}
\end{table}

% ---------------------------------------------------------------------------------------------
\subsection{Relative tests}\label{sec:method-relative}

\paragraph{The anchored relative test.}
The inputs are $\Rnew$ ($X$, $n$ windows), the synthetic sample $G$ ($Y$, $m$ windows) and
$\Rpast$ ($Z$, $k$ windows).
\begin{enumerate}
  \item \emph{Split.} $Y$ and $Z$ are split into training and evaluation halves, of sizes
    $m_t,m_e$ and $k_t,k_e$: by time blocks under the dependent protocol, at random otherwise.
    The block labels are shared, so a block of $G$ (one tenth of \TEST{}) and the block of
    $\Rpast$ with the same label (the corresponding tenth of \TRAIN{}) fall in the same half; in
    E3 this puts about half of each sample in each half. $\Rnew$ is not split: all of it is
    used for training, and none of it for evaluation.
  \item \emph{Training.} One critic $\fhat=\ghat\circ A$, with the architecture and optimiser of
    the \WC{} test and $A$ fitted on all training windows, is trained to separate $\Rnew$
    (positive class) from the pool $Y^{\trn}\cup Z^{\trn}$ (negative class). Inside the pool the
    two parts carry equal total weight, whatever their sizes:
    \begin{equation}\label{eq:anchored-obj}
      \ghat\approx\argmax_g\Big\{\widehat P_{\mathrm{new}}(g\circ A)
      -\tfrac12\QG^{\trn}(g\circ A)-\tfrac12\Ppast^{\trn}(g\circ A)
      -\lgp\,\E\big(\norm{\nabla g(\tilde x)}-1\big)^2\Big\},
    \end{equation}
    where $\widehat P_{\mathrm{new}}$ is the empirical measure of all $\Rnew$ windows and
    $\QG^{\trn},\Ppast^{\trn}$ are those of the two training halves. In practice each mini-batch
    draws pool windows with replacement, with probability $1/(2m_t)$ for each window of
    $Y^{\trn}$ and $1/(2k_t)$ for each window of $Z^{\trn}$, and the penalty is evaluated at
    interpolates between an $\Rnew$ window and a pool window.
  \item \emph{Statistic.} With $\fhat$ frozen,
    \begin{equation}\label{eq:anchored-stat}
      \Danc=\Ppast^{\evl}\fhat-\QG^{\evl}\fhat
      =\frac1{k_e}\sum_{Z_l\in Z^{\evl}}\fhat(Z_l)-\frac1{m_e}\sum_{Y_j\in Y^{\evl}}\fhat(Y_j).
    \end{equation}
\end{enumerate}
The critic scores ``$\Rnew$-likeness''. $\Danc>0$ means that real past windows look more like
the new data than synthetic windows do: $G$ is worse than simply replaying history. The test is
one-sided against $\Danc>0$, and its conditional target is $\delta=\Ppast\fhat-\QG\fhat$, with
null $\delta\le0$ (\cref{def:rel-null}). The name ``anchored'' refers to the role of $\Rnew$.
For one shared critic, the difference of the two absolute discrepancies is
$(\Pnew\fhat-\QG\fhat)-(\Pnew\fhat-\Ppast\fhat)=\Ppast\fhat-\QG\fhat$: the $\Rnew$ terms cancel
exactly. $\Rnew$ anchors what the critic looks for, but is never evaluated, which is why all of
it can be used for training.

\emph{p-values.} For independent windows, the labels of the pooled $G$ and $\Rpast$ evaluation
windows are permuted with $\fhat$ fixed ($B=1000$), and the p-value is \eqref{eq:perm-pvalue}
with $(a,b)=(\fhat(Z^{\evl}),\fhat(Y^{\evl}))$. At the boundary $\Law(G)=\Law(\Rpast)$ this test
is exact for every critic and every law of $\Rnew$ (\cref{prop:anchored}). Given $\Gsig$, the
$G$ and $\Rpast$ evaluation windows are then \iid{} from one law, whatever the training did. The
pool weights affect power only. In the interior of the null, with unequal evaluation sizes, the
plain permutation is only approximately calibrated. The code therefore also reports the
studentised permutation p-value (Welch statistic recomputed under each permutation), which is
asymptotically valid for $\delta\le0$ and still exact at the boundary
\citep{janssen1997studentized,chung2013exact}. For dependent windows the p-value is
\eqref{eq:hac-z} with $(a,b)=(\fhat(Z^{\evl}),\fhat(Y^{\evl}))$, each Driscoll--Kraay variance
computed on its own date ranks. This HAC p-value enters the verdict, and the reported interval
is $\Danc\pm1.96\,\widehat{\mathrm{se}}_{\mathrm{DK}}$.

\emph{Caveats.} (i) $G$ is calibrated on $\Rpast$, so independence between $G$ and the $\Rpast$
windows is an idealisation; for parametric generators each past window moves the fitted
parameters by $O(1/k)$, and we treat the effect as negligible without proof
(\cref{rem:rel-guarantee}). (ii) The stationary bootstrap generator copies blocks of $\Rpast$.
Some synthetic evaluation windows then duplicate windows that the critic saw as pool negatives,
exchangeability fails, and the test is invalid; its verdict is \INVALID{}. E2 part e runs this
case with a naive and with a grouped split (\cref{tab:e2-relative-boundary}). (iii) On real data
$\Rpast$ is not stationary (2020, 2022). The boundary is then the pooled \TRAIN{} window law,
and the time-block split compares the critic's view of $G$ with a random half of the \TRAIN{}
history. A single-seed p-value is not calibrated in this situation. The verdict therefore uses
the seed-combined p-value of \cref{sec:method-decision}, whose power and boundary size are
estimated under a design with the real \TRAIN{} volatility path (design D2 below). (iv) The last \TRAIN{} block
borders the first \TEST{} dates. Under volatility clustering the critic can learn a regime
shared by $\Rnew$ and the end of $\Rpast$, a small effect in the direction of rejection.

\paragraph{The two-critic relative test (secondary).}
All three samples are split, with shared block labels under the dependent protocol. Two critics
are trained separately with the same settings: $\fhat_G$ on $(X^{\trn},Y^{\trn})$ and $\fhat_P$
on $(X^{\trn},Z^{\trn})$. The statistic is the difference of the two held-out discrepancies,
\begin{equation}\label{eq:two-critic-stat}
  \Dtwo=\big(\Pev\fhat_G-\QG^{\evl}\fhat_G\big)-\big(\Pev\fhat_P-\Ppast^{\evl}\fhat_P\big),
\end{equation}
a sum of three evaluation means given $\Gsig$. The p-value comes from a paired bootstrap with
$B=500$ replicates. In each replicate the $\Rnew$ evaluation windows are resampled once and used
for both critics, and the $G$ and $\Rpast$ evaluation windows are resampled independently. The
p-value is centred, $(1+\#\{\Dtwo_{*b}-\Dtwo\ge\Dtwo\})/(1+B)$ \citep{hall1991two}, and the
interval is the percentile interval; under the dependent protocol the three Driscoll--Kraay
variances are added instead. This test is valid conditionally on the two trained critics
(\cref{prop:paired-boot}) but not exact at the boundary $\Law(G)=\Law(\Rpast)$, where its
unconditional size is at least $\alpha$ (\cref{prop:boundary-size}). There, the conditional target $\Dcrit_G-\Dcrit_P$ is the difference
of two exchangeable replicates of a quantity that depends on training. It is symmetric about
zero, with a spread of the order of the evaluation noise, and for fixed critics and growing
evaluation sets the rejection probability tends to $1/2$. In E2 part e, with $n=500$, the excess
was too small to be seen (\cref{sec:sim-relative}). We report the test for comparison only; it
never enters the verdict.

\paragraph{Kernel version.}
For MMD and \SigMMD{} there is no learned witness. The relative test then follows
\citet{bounliphone2016test}: $\TMMD(\Rnew,G)-\TMMD(\Rnew,\Rpast)$ with one kernel (median
heuristic on the three pooled samples). Its normal p-value uses the variance of the first-order
Hoeffding projections, with Driscoll--Kraay variances of the projections under dependence. The
approximation degenerates when all three laws coincide. This version is available in the code
but is not used for any verdict.

% ---------------------------------------------------------------------------------------------
\subsection{The decision rule}\label{sec:method-decision}

\paragraph{Seeds.}
A single anchored p-value depends on the synthetic path and on the randomness of the test (block
assignment, initialisation, mini-batches). Reporting the most favourable of several draws would
be invalid. E3 therefore runs $S=5$ synthetic seeds per generator. Seed $s$ draws a fresh
synthetic path for every ticker and fresh test randomness, and the seeds are shared across
generators, so at a given seed every generator faces the same block assignment and critic
initialisation (common random numbers). With $p_s$ the HAC p-value of the anchored test at seed
$s$, the seeds are combined by
\begin{equation}\label{eq:pmed}
  \pmed=\min\big\{1,\;2\med(p_1,\dots,p_S)\big\},
\end{equation}
which is a valid p-value under arbitrary dependence between the $p_s$
\citep{meinshausen_2009_pvalues}; see \cref{rem:combine-seeds}. With $S=5$ the median is the
third smallest $p_s$.

\paragraph{Power check.}
A non-rejection is evidence of quality only if the test could have detected a deviation of
relevant size. The power check (\cref{def:power-check}) simulates the whole rule, five seeds
included, with the anchored \WC{} test, the dependent protocol and the real sample sizes,
scales, date ranks and split. The pre-specified deviation $\vartheta^*$ multiplies synthetic
volatility by 1.2. There are two reference designs.

\emph{D1 (stationary).} The reference model is the per-ticker \GARCHt{} fitted on
\TRAIN{} (the E3 fits). In each replication, $\Rpast^*$ is a \TRAIN{}-length path started
from stationarity (250 burn-in days), and $\Rnew^*$ is an independent \TEST{}-length path
started from the fitted end-of-\TRAIN{} state. For each seed the alternative $G^*$ is 1.2
times an independent \TEST{}-length draw from the end-of-\TRAIN{} state. The boundary
generator $G^0$ is an independent \TEST{}-length draw started from stationarity, so that
$\Law(G^0)=\Law(\Rpast^*)$ exactly. $G^*$ and $G^0$ are tested with the same test randomness.
The D1 alternative thus differs from the boundary in its start state as well as in its
scale, so its power is not that of the factor 1.2 alone.

\emph{D2 (real volatility path).} Let $v_{ij}$ be the root mean square of the
normalised real \TRAIN{} window $j$ of ticker $i$. Window $(i,j)$ of $\Rpast^*$ is
$v_{ij}\,\xi$ with $\xi\sim N(0,I_L)$, drawn once per replication; this keeps the 2020
and 2022 regimes and their co-movement across tickers. $\Rnew$ is the real \TEST{} sample.
For each seed, the $j'$-th window of ticker $i$ in $G^0$ is $v_{iU}\,\xi'$ with $U$ uniform
over the \TRAIN{} windows of ticker $i$, so $G^0$ has exactly the pooled \TRAIN{} window law
of each ticker; it carries the \TEST{} date ranks. The alternative is $G^*=1.2\,G^0$. The
boundary of D2 is a placebo of the anchored test under the real non-stationarity of
$\Rpast$, with two simplifications. Within a window the returns of $\Rpast^*$ and $G^0$ are Gaussian with a
constant scale, and the noise is independent across tickers and windows; $G^0$ also draws
its scales independently across tickers and windows, so it has no within-window tails or
clustering and no regime common to several tickers. The D2 boundary size is therefore an
estimate for this approximation of the pooled past, not for every generator at the
boundary, in particular not for fat-tailed generators with volatility clustering such as
\GARCHt{}.

Each design runs $R=100$ replications. The estimated power $\hat\beta_{\mathrm{D}}$ is the share
of replications with $\pmed\le\alpha$ under the alternative; the boundary size
$\hat\alpha_{\mathrm{bd},\mathrm{D}}$ is the same share under $G^0$. The power is reported with
its Wilson interval and its one-sided Clopper--Pearson lower bound $\betaL$
\citep{clopper_1934_use}, the boundary size with its Clopper--Pearson interval. The power of the
rule is the smaller of the two, $\hat\beta=\min(\hat\beta_{\mathrm{D1}},\hat\beta_{\mathrm{D2}})$. The designs, the deviation and
the threshold were fixed, and both designs were run, before the full real-data tests; D2 and the
seed rule entered the protocol on the day of those runs, when the real-data tests had been run
only in smoke runs on a six-ticker subset. The check does not involve the generators being
evaluated, but D2 uses the real \TEST{} sample, so its power and boundary size are conditional on
the realised \TEST{} sample and \TRAIN{} volatility path (\cref{tab:e3-power}).

\paragraph{Verdict.}
With $\alpha=0.05$ and $\bmin=0.8$,
\begin{equation}\label{eq:verdict}
  V=\begin{cases}
    \BAD & \text{if } \pmed\le\alpha,\\
    \GOOD & \text{if } \pmed>\alpha \text{ and } \hat\beta\ge\bmin,\\
    \INCONCL & \text{otherwise,}
  \end{cases}
\end{equation}
and $V=\INVALID$ for the stationary bootstrap, whose relative test is invalid (caveat (ii)
above). \Cref{def:verdict} restates the rule formally, and \cref{prop:verdict} states what it
controls. The variant that requires the Clopper--Pearson bounds of both designs to reach $\bmin$
is reported alongside. Each verdict comes with $\pmed$ of the absolute \WC{} test over the same
seeds, the median and range of $\Danc$ over seeds, and the boundary sizes of both designs
(\cref{tab:e3-verdicts}). The combination ``absolute test rejects, relative test does not''
reads ``distinguishable from the future, but not worse than history''. It is a finding, not a
failure of the method.

\paragraph{What the verdict means.}
\BAD{} is a rejection of the anchored null for the trained critics. With independent windows,
the permutation version of the anchored p-value would bound its error rate at the boundary by
$\alpha$ in finite samples (\cref{prop:anchored}(a)), and $\pmed$ would inherit the bound
(\cref{rem:combine-seeds}). The rule as implemented combines Driscoll--Kraay p-values instead,
which are only asymptotically valid (\cref{prop:anchored}(c)), under conditions that the
interleaved split meets only approximately (\cref{sec:method-inference}). The seed combination
is valid only if the single-seed p-values are. In D2 they are not: the single-seed boundary size
is 0.25 [0.17, 0.35] (0.09 [0.04, 0.16] in D1), and the boundary size of the rule is
0.08 [0.04, 0.15] (0.00 [0.00, 0.04] in D1; Clopper--Pearson 95\% intervals,
\cref{tab:e3-power}). On real data the error rate of \BAD{} at the boundary is therefore this
estimated boundary size, which we report rather than assume, and a \BAD{} verdict has a real-data
error interpretation only when the Clopper--Pearson interval of the estimated boundary size covers
$\alpha$ or lies below it (\cref{prop:verdict}(a)), as at the main split. The size-adjusted power
of \cref{def:power-check} separates the discrimination of the rule from this miscalibration
(\cref{tab:manual-theory-calibration}), and \cref{sec:emp-verdicts} also compares each observed
$\pmed$ with the D2 boundary replications.

\GOOD{} means ``no detectable deviation of the pre-specified size and type in designs that
resemble the data''. It never means ``equal to the truth''. Its guarantee rests on a premise
that cannot be verified: the actual deviation must be at least as detectable as a 20\%
volatility error in D1 and D2 (\cref{prop:verdict}). The premise is also one-directional.
$\vartheta^*$ is an upward volatility error; in D2, where $\Rnew$ is the real \TEST{} sample and
\TEST{} was calmer than \TRAIN{}, it moves the generator away from the evaluation period, which
is the direction in which the anchored test is built to reject. The check covers neither
downward volatility errors nor errors in other features, such as tails or clustering, and its
power against a factor 0.8 was not estimated. A downward
error moves $G$ towards \TEST{} and need not violate the anchored null at all, since that null
asks only whether $G$ is further from $\Rnew$ than the past is (\cref{rem:anchored-null}). Under
a non-stationary past, a non-rejection can therefore reflect a generator bias that happens to
point in the direction of the regime change, and a \GOOD{} verdict would credit that bias as
quality. In E3 the rule power in D2 is 0.45, so no generator can be \GOOD{} and this restriction
changes no reported verdict, but it limits what a \GOOD{} verdict means in any application with
a regime change between the two periods. The per-ticker tests of E3 ($L=5$,
calibrated p-values, Benjamini--Hochberg; \cref{sec:emp-perticker,rem:bh}) are descriptive and
do not enter the verdict. \Cref{fig:alg-verdict} collects the steps.

\begin{figure}[tbp]
  \centering
  \fbox{\begin{minipage}{0.95\linewidth}
  \small
  \textbf{Input.} Real \TRAIN{} and \TEST{} panels (100 tickers); a generator class
  $\mathcal M$; $L=20$, $K=10$, $S=5$, $\alpha=0.05$, $\bmin=0.8$, $R=100$.
  \begin{enumerate}
    \setlength{\itemsep}{1pt}
    \item \emph{Power check} (once, for all generators). For $\mathrm D\in\{\mathrm{D1},\mathrm{D2}\}$
      and $r=1,\dots,R$: draw $\Rpast^*$ (and $\Rnew^*$ in D1); for $s=1,\dots,S$ draw $G^*_s$ and
      $G^0_s$ and run step 4 on each; set $p^{\mathrm{alt}}_r$ and $p^{\mathrm{bd}}_r$ by
      \eqref{eq:pmed}. Then $\hat\beta_{\mathrm D}=R^{-1}\sum_r\one\{p^{\mathrm{alt}}_r\le\alpha\}$,
      $\hat\alpha_{\mathrm{bd},\mathrm D}=R^{-1}\sum_r\one\{p^{\mathrm{bd}}_r\le\alpha\}$ and
      $\hat\beta=\min_{\mathrm D}\hat\beta_{\mathrm D}$.
    \item \emph{Fit.} Fit $\mathcal M$ on the \TRAIN{} returns of each ticker.
    \item \emph{Simulate} (for $s=1,\dots,S$). One \TEST{}-length path per ticker with seed $s$,
      from the fitted end-of-\TRAIN{} state where the model has one; divide by $s_i$; cut
      windows; assign \TEST{} date ranks.
    \item \emph{Anchored test} (for each $s$). Time-block split of $G_s$ and $\Rpast$ (shared
      labels); train one \WC{} on all of $\Rnew$ against the pool with weights $\tfrac12,\tfrac12$
      \eqref{eq:anchored-obj}; compute $\Danc_s$ \eqref{eq:anchored-stat} and
      $p_s=1-\Phi(\zhac)$ from \eqref{eq:dk-var}--\eqref{eq:hac-z}.
    \item \emph{Combine.} $\pmed=\min\{1,2\med_sp_s\}$.
    \item \emph{Verdict.} \INVALID{} if $G$ resamples $\Rpast$ (stationary bootstrap); else
      \BAD{} if $\pmed\le\alpha$, \GOOD{} if $\hat\beta\ge\bmin$, \INCONCL{} otherwise.
    \item \emph{Report.} $V$; the Clopper--Pearson variant; $\hat\beta_{\mathrm D}$ and
      $\hat\alpha_{\mathrm{bd},\mathrm D}$; the median and range of $\Danc_s$; $\pmed$ of the
      absolute \WC{} test (time-block split, $\phac$) over the same seeds.
  \end{enumerate}
  \end{minipage}}
  \caption{The decision rule for one generator, as run in E3. Step 1 does not involve the
    generators under evaluation.}
  \label{fig:alg-verdict}
\end{figure}
